\documentclass[11pt]{article}
\usepackage[margin=1in]{geometry}
\usepackage{amsmath,amssymb,amsthm,mathtools}
\usepackage{bm}
\usepackage{enumitem}
\usepackage{hyperref}
 
%% Theorem environments
\newtheorem{theorem}{Theorem}
\newtheorem{lemma}[theorem]{Lemma}
\newtheorem{proposition}[theorem]{Proposition}
\newtheorem{corollary}[theorem]{Corollary}
\newtheorem{remark}[theorem]{Remark}
\theoremstyle{definition}
\newtheorem{definition}[theorem]{Definition}
\newtheorem{assumption}{Assumption}[section]
 
%% Shorthand commands
\newcommand{\Rt}{\mathcal{R}(t)}
\newcommand{\Rtc}{[\mathcal{R}(t)]^c}
\newcommand{\Sig}{\Sigma}
\newcommand{\SigR}{\Sigma_{\Rt}}
\newcommand{\SigRc}{\Sigma_{\Rtc}}
\newcommand{\SigRRc}{\Sigma_{\Rt,\Rtc}}
\newcommand{\SigRcR}{\Sigma_{\Rtc,\Rt}}
\newcommand{\SigSchur}{\Sigma_{\Rtc|\Rt}}
\newcommand{\hSig}{\widehat{\Sigma}}
\newcommand{\bhat}{\hat{\beta}}
\newcommand{\ohat}{\hat{\omega}}
\newcommand{\E}{\mathbb{E}}
\newcommand{\tr}{\mathrm{tr}}
\newcommand{\Cov}{\mathrm{Cov}}
\newcommand{\Var}{\mathrm{Var}}
\newcommand{\op}{\mathrm{op}}
\newcommand{\as}{\mathrm{a.s.}}
\DeclareMathOperator*{\argmin}{arg\,min}
\newcommand{\R}{\mathbb{R}}
\newcommand{\lmax}{\lambda_{\max}}
\newcommand{\lmin}{\lambda_{\min}}
 
\title{Asymptotic Predictive Risk Bound for the\\[4pt]
  Ridge-Regularised Transfer Learning Estimator\\[4pt]
  in High Dimensions}
\author{}
\date{}
 
\begin{document}
\maketitle
 
%% ============================================================
\section{Problem Setup and Notation}\label{sec:setup}
%% ============================================================
 
Consider a target site $t$ and a source site $s$.  For each site
$m\in\{t,s\}$ the response $y^{(m)}\in\R$ and the covariate vector
$x\in\R^p$ satisfy
\[
  y^{(m)} = x^\top \beta^{(m)}+\varepsilon,
  \qquad x\sim N(0,\Sig),\quad
  \varepsilon\sim N(0,\sigma^2),
\]
where $\varepsilon$ is independent of $x$ and $\beta^{(m)}$.
At the target site only the sub-vector $x_{\Rt}$, indexed by the set
$\Rt\subseteq[p]$ with $|\Rt|=p_t$, is observed.  Write
$X\equiv X^{(t)}_{\Rt}\in\R^{n_t\times p_t}$ for the target design
matrix and $Y\equiv Y^{(t)}\in\R^{n_t}$ for the target response.
 
Given an external estimator $\bhat^{(s)}$ of $\beta^{(s)}$, define the
\emph{projection matrix}
$\Phi^{(t)}\in\R^{p_t\times p}$ by
\begin{equation}\label{eq:Phi}
  \bigl[\Phi^{(t)}\bigr]_{\cdot,\,\Rt}=I_{p_t},
  \qquad
  \bigl[\Phi^{(t)}\bigr]_{\cdot,\,\Rtc}
  =\SigR^{-1}\SigRRc,
\end{equation}
and set $\ohat_s^{(t)}=\Phi^{(t)}\bhat^{(s)}$.  The transfer learning
estimator is
\begin{equation}\label{eq:TL}
  \ohat_{\mathrm{TL}}^{(t)}
  :=\argmin_{\omega\in\R^{p_t}}
  \Bigl[\frac{1}{n_t}\|Y-X\omega\|_2^2
        +\lambda_\omega\|\omega-\ohat_s^{(t)}\|_2^2\Bigr],
  \qquad \lambda_\omega>0.
\end{equation}
For an independent test pair
$(\tilde x_{\Rt},\tilde y^{(t)})$ from the target distribution, the
\emph{conditional predictive risk} is
\[
  r^{\mathrm{pred}}\!\bigl(\ohat_{\mathrm{TL}}^{(t)}\bigr)
  =\E\!\Bigl[
     \bigl(\tilde y^{(t)}
           -\tilde x_{\Rt}^\top\ohat_{\mathrm{TL}}^{(t)}\bigr)^2
     \;\Big|\; X\Bigr].
\]
 
%% ============================================================
\section{Assumptions}\label{sec:assumptions}
%% ============================================================
 
\begin{assumption}[High-Dimensional Proportional Limit]\label{as:dim}
  As $n_t\to\infty$, $p/n_t\to\gamma\in(0,\infty)$ and
  $p_t/n_t\to\eta\in(0,\infty)$.  In particular,
  $p_t/p\to\eta/\gamma\in(0,1]$.
\end{assumption}
 
\begin{assumption}[Random Regression Coefficients]\label{as:beta}
  For each $m\in\{s,t\}$,
  $\E[\beta^{(m)}]=0$ and
  $\Cov(\beta^{(m)})=(\alpha_m^2/p)\,I_p$ with $\alpha_m\neq 0$.
  The cross-covariance satisfies
  $\Cov(\beta^{(t)},\beta^{(s)})=(\rho\alpha_t\alpha_s/p)\,I_p$ with
  $\rho\in(0,1]$.
  Moreover, $\beta^{(t)}$ and $\beta^{(s)}$ are independent of the
  target design~$X$, the test point~$(\tilde x,\tilde y^{(t)})$, and the
  target noise.
\end{assumption}
 
\begin{assumption}[Source Estimation Error]\label{as:source}
  The source estimator $\bhat^{(s)}$ is independent of the target data
  $(X,Y)$ and the test observation, and satisfies
  \[
    \lmax\!\Bigl(
      \E\bigl[(\bhat^{(s)}-\beta^{(s)})
              (\bhat^{(s)}-\beta^{(s)})^\top\bigr]
    \Bigr)\le \frac{C^U}{p}
  \]
  uniformly as $(n_t,p)\to\infty$, for a constant $C^U>0$.
\end{assumption}
 
\begin{assumption}[Population Covariance Regularity]\label{as:Sigma}
  The eigenvalues of $\Sig$ are uniformly bounded:
  \[
    0<c_\ell\le\lmin(\Sig)\le\lmax(\Sig)\le c_u<\infty
  \]
  for constants $c_\ell,c_u$ independent of $p$.  In particular,
  $\SigR$ inherits the same bounds.  The empirical spectral distribution
  of $\SigR$,
  \[
    H_{p_t}(\cdot)
    =\frac{1}{p_t}\sum_{j=1}^{p_t}\delta_{\tau_j}(\cdot),
    \qquad
    \tau_1\ge\cdots\ge\tau_{p_t}>0
    \text{ eigenvalues of }\SigR,
  \]
  converges weakly to a compactly supported probability measure $H$ on
  $(0,\infty)$.
\end{assumption}
 
\begin{assumption}[Bounded Projection Operator]\label{as:Phi}
  The matrix $\Gamma:=\Phi^{(t)}(\Phi^{(t)})^\top$ satisfies
  $\lmax(\Gamma)\le c_\Phi$ uniformly, for a constant $c_\Phi>0$.
\end{assumption}
 
\begin{remark}
Assumption~\ref{as:Phi} is implied by Assumption~\ref{as:Sigma}.
Indeed, $\Phi^{(t)}=\SigR^{-1}\Sig_{\Rt,[p]}$ where
$\Sig_{\Rt,[p]}$ is the $p_t\times p$ sub-matrix of~$\Sig$.  Then
$\Gamma=\SigR^{-1}(\Sig^2)_{\Rt,\Rt}\SigR^{-1}$ and hence
$\lmax(\Gamma)\le \|\SigR^{-1}\|_{\op}^2\|\Sig\|_{\op}^2
\le c_u^2/c_\ell^2$.
\end{remark}
 
%% ============================================================
\section{Statement of the Main Result}\label{sec:main}
%% ============================================================
 
\begin{definition}[Marchenko--Pastur companion Stieltjes transform]\label{def:MP}
For $\lambda>0$, let $v(\lambda)>0$ be the unique positive solution of
the fixed-point equation
\begin{equation}\label{eq:FP}
  \frac{1}{v(\lambda)}
  =\lambda+\eta\!\int\!\frac{\tau}{1+\tau\,v(\lambda)}\,dH(\tau),
\end{equation}
where $H$ is the limiting spectral distribution of $\SigR$
(Assumption~\ref{as:Sigma}) and
$\eta=\lim p_t/n_t$.
\end{definition}
 
\begin{definition}[Asymptotic trace functionals]\label{def:traces}
Define
\begin{align}
  b_0(\lambda)
  &:=\int\!\frac{\tau}{(1+\tau\,v(\lambda))^2}\,dH(\tau)
  \cdot
  \frac{v(\lambda)^2}{1-\eta\!\displaystyle\int
    \frac{\tau^2\,v(\lambda)^2}{(1+\tau\,v(\lambda))^2}\,dH(\tau)},
  \label{eq:b0}
  \\[6pt]
  v_0(\lambda)
  &:=
  \eta\!\int\!\frac{\tau}{(1+\tau\,v(\lambda))^2}\,dH(\tau)
  \cdot
  \frac{1}{1-\eta\!\displaystyle\int
    \frac{\tau^2\,v(\lambda)^2}{(1+\tau\,v(\lambda))^2}\,dH(\tau)}.
  \label{eq:v0}
\end{align}
\end{definition}
 
\begin{theorem}[Asymptotic Predictive Risk Upper Bound]\label{thm:main}
Under Assumptions
\ref{as:dim}--\ref{as:Phi}, for every fixed
$\lambda_\omega>0$ the conditional predictive risk of the estimator
\eqref{eq:TL} satisfies, almost surely as $(n_t,p)\to\infty$,
\begin{equation}\label{eq:risk_bound}
  \boxed{
  \limsup_{n_t,p\to\infty}
  \E\!\left[
  r^{\mathrm{pred}}\!\bigl(\ohat_{\mathrm{TL}}^{(t)}\bigr)
  \right]
  \;\le\;
  \mathcal R_H\!\left(\lambda_\omega;B_\delta,\sigma_0^2,\eta\right),
  }
\end{equation}
where
\begin{itemize}[nosep]
\item $\sigma_0^2
  =\dfrac{\alpha_t^2}{p}\,\tr(\SigSchur)+\sigma^2$
  is the irreducible prediction variance
  (which converges to a finite limit under
  Assumption~\ref{as:Sigma}),
\item $C_\delta=2\bigl(h^2+C^U\bigr)$ with
  $h^2=\alpha_t^2+\alpha_s^2-2\rho\,\alpha_t\alpha_s$
  being the transfer discrepancy,
\item $B_\delta=(\eta/\gamma)c_\Phi C_\delta$ is the effective
  signal-strength parameter entering the ordinary ridge-risk
  functional,
\item $c_\Phi\le c_u^2/c_\ell^2$ is the uniform bound on
  $\lmax(\Gamma)$,
\item $\mathcal R_H(\lambda;a^2,\tau^2,\eta)$ is the deterministic
  equivalent predictive-risk functional for ridge regression with
  population covariance spectral law $H$, signal strength $a^2$, noise
  variance $\tau^2$, and aspect ratio $\eta$.
\end{itemize}
\end{theorem}
 
\begin{corollary}[Optimal Regularisation Parameter]\label{cor:opt}
The optimally tuned upper bound can be written without expanding the
deterministic-equivalent ridge-risk functional as
\begin{equation}\label{eq:opt_lambda}
  \inf_{\lambda_\omega>0}
  \limsup_{n_t,p\to\infty}
  \E\!\left[
  r^{\mathrm{pred}}\!\bigl(\ohat_{\mathrm{TL}}^{(t)}\bigr)
  \right]
  \;\le\;
  \inf_{\lambda_\omega>0}
  \mathcal R_H\!\left(\lambda_\omega;B_\delta,\sigma_0^2,\eta\right).
\end{equation}
In the isotropic case $\Sig=I_p$, this is minimized at
\begin{equation}\label{eq:opt_iso}
  \lambda_\omega^*
  =\frac{\gamma\,\sigma_0^2}{c_\Phi C_\delta}
  =\frac{\gamma\,\sigma_0^2}{2(h^2+C^U)}
  \qquad(c_\Phi=1),
\end{equation}
and the corresponding unexpanded conclusion is
\begin{equation}\label{eq:opt_iso_unexpanded}
  \boxed{
  \inf_{\lambda_\omega>0}
  \limsup_{n_t,p\to\infty}
  \E\!\left[
  r^{\mathrm{pred}}\!\bigl(\ohat_{\mathrm{TL}}^{(t)}\bigr)
  \right]
  \;\le\;
  \mathcal R_{\delta_1}\!\left(
  \lambda_\omega^*;(\eta/\gamma)C_\delta,\sigma_0^2,\eta
  \right).
  }
\end{equation}
\end{corollary}
 
%% ============================================================
\section{Proof of Theorem~\ref{thm:main}}\label{sec:proof}
%% ============================================================
 
The proof proceeds in five steps:
(i) closed-form solution of the estimator;
(ii) oracle target coefficient and noise decomposition;
(iii) predictive risk decomposition into bias, variance, and irreducible
error;
(iv) bounding the bias via the transfer discrepancy;
(v) asymptotic evaluation of trace functionals using the
Marchenko--Pastur deterministic equivalents of
Dobriban and Wager~(2018).
 
\subsection{Step 1: Closed-Form Solution}\label{sec:step1}
 
The objective in \eqref{eq:TL} is strongly convex.  Setting the
gradient to zero gives
\begin{equation}\label{eq:closed}
  \ohat_{\mathrm{TL}}^{(t)}
  =\bigl(\hSig+\lambda_\omega I_{p_t}\bigr)^{-1}
  \Bigl(\frac{1}{n_t}X^\top Y+\lambda_\omega\,\ohat_s^{(t)}\Bigr),
\end{equation}
where $\hSig=X^\top X/n_t\in\R^{p_t\times p_t}$ is the sample Gram
matrix.  We abbreviate $A(\lambda):=(\hSig+\lambda I_{p_t})^{-1}$ and
write $A\equiv A(\lambda_\omega)$.
 
\subsection{Step 2: Oracle Target Coefficient}\label{sec:step2}
 
Because only $x_{\Rt}$ is observed at the target site, the best
(minimum mean-squared-error) linear predictor of $y^{(t)}$ given
$x_{\Rt}$ (conditional on~$\beta^{(t)}$) uses the coefficient vector
\begin{equation}\label{eq:omega_star}
  \omega^*:=\SigR^{-1}\Sig_{\Rt,[p]}\,\beta^{(t)}
  =\Phi^{(t)}\beta^{(t)},
\end{equation}
where $\Sig_{\Rt,[p]}=\E[x_{\Rt}\,x^\top]$ is the $p_t\times p$
cross-covariance.  This follows from the well-known formula for the
conditional expectation of a jointly Gaussian vector:
$\E[x_{\Rtc}\mid x_{\Rt}]
=\SigRcR\SigR^{-1}x_{\Rt}$, so that
\[
  \E[y^{(t)}\mid x_{\Rt},\beta^{(t)}]
  =x_{\Rt}^\top
  \bigl(\beta^{(t)}_{\Rt}
  +\SigR^{-1}\SigRRc\,\beta^{(t)}_{\Rtc}\bigr)
  =x_{\Rt}^\top\omega^*.
\]
Define the \emph{target residual vector}
$\xi\in\R^{n_t}$ entry-wise by
\begin{equation}\label{eq:xi}
  \xi_i := y_i^{(t)}-x_{i,\Rt}^\top\omega^*
  =\bigl(x_{i,\Rtc}-\SigRcR\SigR^{-1}x_{i,\Rt}\bigr)^\top
    \beta^{(t)}_{\Rtc}+\varepsilon_i,
  \qquad i=1,\dots,n_t.
\end{equation}
The vector
$x_{i,\Rtc|\Rt}:=x_{i,\Rtc}-\SigRcR\SigR^{-1}x_{i,\Rt}$
is the conditional residual, which by Gaussianity is independent of
$x_{i,\Rt}$ and has covariance
$\SigSchur=\SigRc-\SigRcR\SigR^{-1}\SigRRc$
(the Schur complement).  Therefore, conditional on $\beta^{(t)}$,
the $\xi_i$ are independent of $X$ with
\begin{equation}\label{eq:sigxi}
  \E[\xi_i\mid\beta^{(t)}]=0,
  \qquad
  \Var[\xi_i\mid\beta^{(t)}]
  =(\beta^{(t)}_{\Rtc})^\top\SigSchur\,\beta^{(t)}_{\Rtc}
  +\sigma^2.
\end{equation}
In matrix form, $Y=X\omega^*+\xi$.
 
\subsection{Step 3: Risk Decomposition}\label{sec:step3}
 
For the test point we write similarly
$\tilde y^{(t)}=\tilde x_{\Rt}^\top\omega^*+\tilde\xi$
with $\tilde\xi$ independent of $\tilde x_{\Rt}$ (conditional on
$\beta^{(t)}$) and having the same variance
\eqref{eq:sigxi}.
 
\begin{lemma}[Bias--Variance--Noise Decomposition]\label{lem:decomp}
\begin{equation}\label{eq:decomp}
  r^{\mathrm{pred}}\!\bigl(\ohat_{\mathrm{TL}}^{(t)}\bigr)
  =\underbrace{\E\!\bigl[
      (\ohat_{\mathrm{TL}}^{(t)}-\omega^*)^\top
      \SigR\,
      (\ohat_{\mathrm{TL}}^{(t)}-\omega^*)
    \;\big|\;X\bigr]}_{\displaystyle=:R_{\mathrm{est}}(X)}
  \;+\;\sigma_0^2,
\end{equation}
where
\begin{equation}\label{eq:sig0}
  \sigma_0^2
  :=\E[\tilde\xi^{\,2}]
  =\frac{\alpha_t^2}{p}\,\tr(\SigSchur)+\sigma^2.
\end{equation}
\end{lemma}
 
\begin{proof}
Since $\tilde x_{\Rt}\sim N(0,\SigR)$ is independent of
$X$, $Y$, and $\tilde\xi$, and since $\tilde\xi$ is independent of
$\tilde x_{\Rt}$ (conditional on~$\beta^{(t)}$), the cross-term
vanishes:
\begin{align*}
  r^{\mathrm{pred}}
  &=\E\!\bigl[
    \bigl(\tilde x_{\Rt}^\top(\omega^*-\ohat_{\mathrm{TL}}^{(t)})
          +\tilde\xi\bigr)^2
    \;\big|\;X\bigr]\\
  &=\E\!\bigl[
    (\omega^*-\ohat_{\mathrm{TL}}^{(t)})^\top
    \tilde x_{\Rt}\tilde x_{\Rt}^\top
    (\omega^*-\ohat_{\mathrm{TL}}^{(t)})
    \;\big|\;X\bigr]
    +\E[\tilde\xi^{\,2}]
    +2\,\E\!\bigl[
      \tilde\xi\,
      \tilde x_{\Rt}^\top(\omega^*-\ohat_{\mathrm{TL}}^{(t)})
    \;\big|\;X\bigr].
\end{align*}
The cross-term is zero because, upon first conditioning on
$\beta^{(t)}$ and $\ohat_{\mathrm{TL}}^{(t)}$,
$\E[\tilde\xi\mid\beta^{(t)}]=0$ and
$\tilde x_{\Rt}$ is independent of both $\tilde\xi$ and
$\ohat_{\mathrm{TL}}^{(t)}$ (given~$X$).  Taking the outer expectation
over $\tilde x_{\Rt}$ yields $\SigR$.
 
For~\eqref{eq:sig0}, use Assumption~\ref{as:beta}:
$\E[(\beta^{(t)}_{\Rtc})^\top\SigSchur\,\beta^{(t)}_{\Rtc}]
=(\alpha_t^2/p)\tr(\SigSchur)$.
\end{proof}
 
It remains to bound $R_{\mathrm{est}}(X)$.  Substituting $Y=X\omega^*+\xi$
into~\eqref{eq:closed} and using the identity
$A\hSig=I-\lambda_\omega A$ gives
\begin{align}
  \ohat_{\mathrm{TL}}^{(t)}-\omega^*
  &=A\!\left(\hSig\,\omega^*
    +\frac{X^\top\xi}{n_t}
    +\lambda_\omega\ohat_s^{(t)}\right)-\omega^*
  \notag\\
  &=(I-\lambda_\omega A)\omega^*
    +A\frac{X^\top\xi}{n_t}
    +\lambda_\omega A\ohat_s^{(t)}
    -\omega^*
  \notag\\
  &=-\lambda_\omega A\bigl(\omega^*-\ohat_s^{(t)}\bigr)
    +A\frac{X^\top\xi}{n_t}.
  \label{eq:delta}
\end{align}
Define the two components:
\begin{equation}\label{eq:B1B2}
  B_1:=-\lambda_\omega A\bigl(\omega^*-\ohat_s^{(t)}\bigr),
  \qquad
  B_2:=A\frac{X^\top\xi}{n_t}.
\end{equation}
 
\begin{lemma}[Vanishing Cross-Term]\label{lem:cross}
$\E[B_1^\top\SigR\,B_2\mid X]=0$.
\end{lemma}
 
\begin{proof}
Condition first on $\beta^{(t)},\beta^{(s)},\bhat^{(s)}$ and~$X$.
Then $B_1$ is measurable with respect to this $\sigma$-algebra,
while $B_2=AX^\top\xi/n_t$ has conditional expectation zero because
$\E[\xi\mid X,\beta^{(t)}]=0$ (from~\eqref{eq:sigxi}; the
entries $\xi_i$ have zero mean and are independent of~$X$ given
$\beta^{(t)}$).  Since
$\E[B_1^\top\SigR\,B_2\mid X]
=\E\bigl[B_1^\top\SigR\,
\E[B_2\mid X,\beta^{(t)},\bhat^{(s)}]\;\big|\;X\bigr]=0$.
\end{proof}
 
Therefore
\begin{equation}\label{eq:Rest}
  R_{\mathrm{est}}(X)
  =\underbrace{\E[B_1^\top\SigR\,B_1\mid X]}_{
    \displaystyle\text{Bias}^2}
  +\underbrace{\E[B_2^\top\SigR\,B_2\mid X]}_{
    \displaystyle\text{Variance}}.
\end{equation}
 
\subsection{Step 4: Bounding the Bias}\label{sec:step4}
 
\paragraph{4a.  Transfer discrepancy matrix.}
Since $\omega^*=\Phi^{(t)}\beta^{(t)}$ and
$\ohat_s^{(t)}=\Phi^{(t)}\bhat^{(s)}$,
\[
  \omega^*-\ohat_s^{(t)}
  =\Phi^{(t)}\bigl(\beta^{(t)}-\bhat^{(s)}\bigr)
  =:\Phi^{(t)}\delta_\beta.
\]
 
\begin{lemma}[Transfer Discrepancy Bound]\label{lem:Cdelta}
Under Assumptions~\ref{as:beta} and~\ref{as:source},
\[
  \E[\delta_\beta\,\delta_\beta^\top]
  \preceq \frac{C_\delta}{p}\,I_p,
  \qquad
  C_\delta:=2\bigl(h^2+C^U\bigr),
  \quad
  h^2:=\alpha_t^2+\alpha_s^2-2\rho\,\alpha_t\alpha_s.
\]
\end{lemma}
 
\begin{proof}
Decompose
$\delta_\beta=\underbrace{(\beta^{(t)}-\beta^{(s)})}_{=:u}
+\underbrace{(\beta^{(s)}-\bhat^{(s)})}_{=:w}$.
By the elementary matrix inequality
$(u+w)(u+w)^\top\preceq 2\,uu^\top+2\,ww^\top$ (which follows from
$(u-w)(u-w)^\top\succeq 0$),
\[
  \E[\delta_\beta\delta_\beta^\top]
  \preceq
  2\,\E[uu^\top]+2\,\E[ww^\top].
\]
From Assumption~\ref{as:beta},
\[
  \E[uu^\top]
  =\Cov(\beta^{(t)})+\Cov(\beta^{(s)})
  -\Cov(\beta^{(t)},\beta^{(s)})
  -\Cov(\beta^{(s)},\beta^{(t)})
  =\frac{\alpha_t^2+\alpha_s^2-2\rho\alpha_t\alpha_s}{p}\,I_p
  =\frac{h^2}{p}\,I_p.
\]
From Assumption~\ref{as:source},
$\E[ww^\top]\preceq (C^U/p)\,I_p$.  Combining yields the result.
\end{proof}
 
\paragraph{4b.  Bias evaluation.}
 
\begin{align}
  \text{Bias}^2
  &=\lambda_\omega^2\,
    \E\!\bigl[\delta_\beta^\top(\Phi^{(t)})^\top
    A\SigR A\,\Phi^{(t)}\delta_\beta
    \;\big|\;X\bigr]
  \notag\\
  &=\lambda_\omega^2\,
    \tr\!\Bigl(A\SigR A\,
    \Phi^{(t)}\E[\delta_\beta\delta_\beta^\top](\Phi^{(t)})^\top
    \Bigr)
  \label{eq:bias_trace}
\end{align}
since $\delta_\beta$ is independent of $X$
(Assumptions~\ref{as:beta},~\ref{as:source}).
Applying Lemma~\ref{lem:Cdelta} and
$\Gamma=\Phi^{(t)}(\Phi^{(t)})^\top$,
\begin{align}
  \text{Bias}^2
  &\le
  \lambda_\omega^2\,\frac{C_\delta}{p}\,
  \tr\bigl(A\SigR A\,\Gamma\bigr)
  \notag\\
  &\le
  \lambda_\omega^2\,\frac{C_\delta}{p}\,c_\Phi\,
  \tr\bigl(A^2\SigR\bigr),
  \label{eq:bias_bound}
\end{align}
where the second inequality uses
$\Gamma\preceq c_\Phi I_{p_t}$ (Assumption~\ref{as:Phi})
and the cyclic property of the trace
($\tr(A\SigR A\,\Gamma)\le c_\Phi\,\tr(A\SigR A)
=c_\Phi\,\tr(A^2\SigR)$).
 
\subsection{Step 5: Variance Computation}\label{sec:step5}
 
Conditional on $\beta^{(t)}$ and $X$, the entries $\xi_i$ are
independent with mean~$0$ and common variance
$\sigma_\xi^2(\beta^{(t)})$
from~\eqref{eq:sigxi}.  Therefore
\begin{align}
  \E[B_2^\top\SigR\,B_2\mid X,\beta^{(t)}]
  &=\frac{1}{n_t^2}\,
    \E[\xi^\top X A\SigR A X^\top\xi\mid X,\beta^{(t)}]
  \notag\\
  &=\frac{\sigma_\xi^2(\beta^{(t)})}{n_t^2}\,
    \tr\!\bigl(X A\SigR A X^\top\bigr)
  \notag\\
  &=\frac{\sigma_\xi^2(\beta^{(t)})}{n_t}\,
    \tr\!\bigl(A\SigR A\hSig\bigr),
  \label{eq:var_cond}
\end{align}
where we used $\E[\xi\xi^\top\mid X,\beta^{(t)}]
=\sigma_\xi^2(\beta^{(t)})\,I_{n_t}$ and the cyclic property.
Taking the expectation over $\beta^{(t)}$ (using~\eqref{eq:sig0}):
\begin{equation}\label{eq:var}
  \text{Variance}
  =\frac{\sigma_0^2}{n_t}\,
  \tr\!\bigl(A^2\SigR\hSig\bigr).
\end{equation}
 
\subsection{Step 6: Asymptotic Evaluation via Marchenko--Pastur Theory}
\label{sec:step6}
 
We now appeal to the deterministic-equivalent theory for sample
covariance matrices (see Dobriban and Wager, 2018, Theorem~1, and the
underlying results of Marchenko and Pastur, 1967; Silverstein and Bai,
1995).
 
Write $X=Z\SigR^{1/2}$ where $Z\in\R^{n_t\times p_t}$ has i.i.d.\
$N(0,1)$ entries.  Then
$\hSig=\SigR^{1/2}(Z^\top Z/n_t)\SigR^{1/2}$.
 
\begin{lemma}[Deterministic Equivalents]\label{lem:det_equiv}
Under Assumptions~\ref{as:dim} and~\ref{as:Sigma}, for each fixed
$\lambda>0$, the following hold almost surely as $n_t\to\infty$:
\begin{align}
  \frac{1}{p_t}\,\tr\!\bigl[(\hSig+\lambda I)^{-1}\bigr]
  &\;\longrightarrow\; v(\lambda),
  \label{eq:DE1}\\[4pt]
  \frac{1}{p_t}\,\tr\!\bigl[(\hSig+\lambda I)^{-2}\SigR\bigr]
  &\;\longrightarrow\; b_0(\lambda),
  \label{eq:DE2}\\[4pt]
  \frac{1}{n_t}\,\tr\!\bigl[(\hSig+\lambda I)^{-2}\SigR\hSig\bigr]
  &\;\longrightarrow\; v_0(\lambda),
  \label{eq:DE3}
\end{align}
where $v(\lambda)$ is given by Definition~\ref{def:MP},
and $b_0(\lambda)$, $v_0(\lambda)$ are given in
Definition~\ref{def:traces}.
\end{lemma}
 
\begin{proof}
\eqref{eq:DE1} is the classical companion Stieltjes transform convergence
for sample covariance matrices with population spectral distribution~$H$
and aspect ratio~$\eta$ (Silverstein and Bai, 1995; Dobriban and Wager,
2018, Theorem~1).
 
For \eqref{eq:DE2}, we use the resolvent identity and differentiation.
Define $g(\lambda):=\frac{1}{p_t}\tr[(\hSig+\lambda I)^{-1}\SigR]$.
By the deterministic equivalent for general linear functionals of the
resolvent (Dobriban and Wager, 2018, Eq.~(7)), almost surely
\[
  g(\lambda)
  \;\longrightarrow\;
  \int\frac{\tau}{1+\tau\,v(\lambda)}\,dH(\tau).
\]
Since $\frac{d}{d\lambda}(\hSig+\lambda I)^{-1}
=-(\hSig+\lambda I)^{-2}$,
\[
  \frac{1}{p_t}\tr\bigl[(\hSig+\lambda I)^{-2}\SigR\bigr]
  =-\frac{d}{d\lambda}\,g(\lambda).
\]
By the uniform convergence of the Stieltjes transform on compact subsets
of $(0,\infty)$ and the fact that the limiting function is analytic,
the derivative converges as well.  Differentiating the limit:
\begin{align*}
  -\frac{d}{d\lambda}\int\frac{\tau}{1+\tau v(\lambda)}\,dH(\tau)
  &=\int\frac{\tau^2 v'(\lambda)}{(1+\tau v(\lambda))^2}\,dH(\tau).
\end{align*}
To compute $v'(\lambda)$, differentiate the fixed-point
equation~\eqref{eq:FP}.  Writing
$F(v,\lambda):=v^{-1}-\lambda
-\eta\int\tau/(1+\tau v)\,dH(\tau)=0$,
by the implicit function theorem,
\[
  v'(\lambda)
  =-\frac{\partial F/\partial\lambda}
         {\partial F/\partial v}
  =-\frac{-1}{-v^{-2}+\eta\int\frac{\tau^2}{(1+\tau v)^2}dH(\tau)}
  =\frac{-v^2}{1-\eta v^2\int\frac{\tau^2}{(1+\tau v)^2}dH(\tau)}.
\]
Substituting and simplifying yields
\eqref{eq:DE2} with $b_0(\lambda)$ as stated in \eqref{eq:b0}.
 
For \eqref{eq:DE3}, use the algebraic identity
$A^2\SigR\hSig = A\SigR - \lambda A^2\SigR$
(from $A\hSig = I - \lambda A$, hence
$A^2\SigR\hSig = A^2(\hSig+\lambda I-\lambda I)\SigR
= A\SigR - \lambda A^2\SigR$).  Therefore
\[
  \frac{1}{n_t}\tr(A^2\SigR\hSig)
  =\frac{p_t}{n_t}\left[
    \frac{1}{p_t}\tr(A\SigR)
    -\lambda\frac{1}{p_t}\tr(A^2\SigR)
  \right]
  \;\longrightarrow\;
  \eta\!\left[
    \int\frac{\tau}{1+\tau v}\,dH(\tau)
    -\lambda\,b_0(\lambda)
  \right]
  =:v_0(\lambda).
\]
Substituting the expression for $b_0(\lambda)$ and simplifying yields
the formula in~\eqref{eq:v0}.
\end{proof}
 
\subsection{Step 7: Assembling the Bound}\label{sec:step7}
 
Combining \eqref{eq:decomp}, \eqref{eq:Rest}, \eqref{eq:bias_bound},
and~\eqref{eq:var}:
\begin{equation}\label{eq:finite_bound}
  r^{\mathrm{pred}}
  \le
  \sigma_0^2
  +\sigma_0^2\cdot\frac{1}{n_t}\tr\bigl(A^2\SigR\hSig\bigr)
  +\lambda_\omega^2\frac{c_\Phi C_\delta}{p}
    \tr\bigl(A^2\SigR\bigr).
\end{equation}
Now apply the deterministic equivalents of Lemma~\ref{lem:det_equiv}:
\begin{align*}
  \frac{1}{n_t}\tr(A^2\SigR\hSig)
  &\;\longrightarrow\; v_0(\lambda_\omega),\\[4pt]
  \frac{1}{p}\tr(A^2\SigR)
  &=\frac{p_t}{p}\cdot
    \frac{1}{p_t}\tr(A^2\SigR)
  \;\longrightarrow\;
  \frac{\eta}{\gamma}\,b_0(\lambda_\omega).
\end{align*}
Substituting into~\eqref{eq:finite_bound} yields the asymptotic
bound~\eqref{eq:risk_bound}:
\[
  r^{\mathrm{pred}}\!\bigl(\ohat_{\mathrm{TL}}^{(t)}\bigr)
  \;\le\;
  \sigma_0^2\bigl(1+v_0(\lambda_\omega)\bigr)
  +\lambda_\omega^2\,\frac{\eta}{\gamma}\,c_\Phi\,C_\delta\;
  b_0(\lambda_\omega)
  +o(1)\qquad\text{a.s.}
\]
This completes the proof of Theorem~\ref{thm:main}.\qed
 
%% ============================================================
\section{Proof of Corollary~\ref{cor:opt}}
\label{sec:proof_opt}
%% ============================================================
 
Theorem~\ref{thm:main} reduces the transfer-learning risk bound to the
ordinary deterministic-equivalent ridge-risk functional
$\mathcal R_H(\lambda_\omega;B_\delta,\sigma_0^2,\eta)$, with
$B_\delta=(\eta/\gamma)c_\Phi C_\delta$.  Taking the infimum over
$\lambda_\omega>0$ immediately gives \eqref{eq:opt_lambda}.

In the isotropic case $H=\delta_1$ and $c_\Phi=1$.  The classical
ridge-risk optimizer has the signal-to-noise form
$\lambda_\omega^*=\eta\sigma_0^2/B_\delta$, which equals
$\gamma\sigma_0^2/C_\delta$ because
$B_\delta=(\eta/\gamma)C_\delta$.  Substituting the coarse discrepancy
constant $C_\delta=2(h^2+C^U)$ gives \eqref{eq:opt_iso}, and
\eqref{eq:opt_iso_unexpanded} is just the same optimized bound written
without expanding $\mathcal R_{\delta_1}$ into its trace formula.\qed
 
%% ============================================================
\section{Discussion}\label{sec:discussion}
%% ============================================================
 
\paragraph{Interpretation.}
The bound~\eqref{eq:risk_bound} decomposes the predictive risk into
three interpretable components:
\begin{enumerate}[label=(\roman*),nosep]
\item \textbf{Irreducible error} $\sigma_0^2$: the cost of observing
  only $p_t$ of the $p$ features, plus the intrinsic noise $\sigma^2$.
  This term vanishes if all features are observed ($p_t=p$, $\Sig=I$)
  and $\sigma^2=0$.
\item \textbf{Variance}
  $\sigma_0^2\,v_0(\lambda_\omega)$: the estimation variance of the
  ridge estimator, amplified by the effective noise $\sigma_0^2$ and
  modulated by the high-dimensional trace functional $v_0$.  This term
  \emph{decreases} in $\lambda_\omega$.
\item \textbf{Bias}
  $\lambda_\omega^2(\eta/\gamma)\,c_\Phi C_\delta\,b_0(\lambda_\omega)$:
  the regularisation bias from shrinking toward the (imperfect) source
  estimate.  This term \emph{increases} in $\lambda_\omega$.  The
  transfer discrepancy $h^2=\alpha_t^2+\alpha_s^2-2\rho\alpha_t\alpha_s$
  measures the distance between the source and target tasks, and
  $C^U$ measures the estimation quality of the source.
\end{enumerate}
 
\paragraph{When does transfer help?}
Transfer learning improves over the target-only ridge estimator
(which corresponds to $\lambda_\omega\to\infty$ or, equivalently,
ordinary ridge with no source information) when $C_\delta$ is small
relative to $\sigma_0^2$, i.e., when the source-target discrepancy
$h^2$ and the source estimation error $C^U$ are jointly small compared
to the target noise level.
 
\paragraph{Connection to Dobriban--Wager and Gu--Han--Duan.}
Our proof strategy follows the general blueprint of
Dobriban and Wager~(2018) for the asymptotic trace computations
(Lemma~\ref{lem:det_equiv}) and the
bias-bounding strategy of Gu, Han, and Duan~(2025) for the
$L_2$-penalised transfer learning estimator
(Steps~3--4, especially the use of the elementary matrix inequality
$(u+w)(u+w)^\top\preceq 2uu^\top+2ww^\top$ and the spectral bound on
the projection matrix~$\Gamma$).
 
\bigskip
\noindent\textbf{References.}
 
\medskip\noindent
\hangindent=1.5em
Dobriban, E.\ and Wager, S.\ (2018).
High-dimensional asymptotics of prediction: Ridge regression and
classification.
\textit{The Annals of Statistics}, 46(1), 247--279.
 
\medskip\noindent
\hangindent=1.5em
Gu, Y., Han, R., and Duan, R.\ (2025).
Robust angle-based transfer learning in high dimensions.
\textit{Preprint}, arXiv:2410.12235.
 
\medskip\noindent
\hangindent=1.5em
Marchenko, V.~A.\ and Pastur, L.~A.\ (1967).
Distribution of eigenvalues of some sets of random matrices.
\textit{Matematicheskii Sbornik}, 114(4), 507--536.
 
\medskip\noindent
\hangindent=1.5em
Silverstein, J.~W.\ and Bai, Z.~D.\ (1995).
On the empirical distribution of eigenvalues of a class of large
dimensional random matrices.
\textit{Journal of Multivariate Analysis}, 54(2), 175--192.
 
\end{document}
